% ═══════════════════════════════════════════════════════════════════════════
%  DRAFT — expansion item #18: "Threats to validity"
%  Target: Chapter 5, after "Non-Topological Context and Limitations" (\S5.2)
%  or as a restructured replacement for it. The content is drawn from the
%  paper's existing Limitations paragraph (\S5.2, tex lines 1751-1781) and the
%  methodological caveats scattered through \S4.1-\S4.3 and \S5.2, reorganised
%  into the standard construct / internal / external validity categories.
%
%  STATUS: DRAFT for review. Not yet \input into dissertation.tex.
%  Each threat: description -> severity -> mitigation (done / planned
%  expansion). Cross-references use the expansion numbering of
%  EXPANSION_PLAN.md.
% ═══════════════════════════════════════════════════════════════════════════

\section{Threats to Validity}
\label{sec:threats}

We structure the limitations of \S5.2 into the three
standard validity categories. The purpose is not to multiply caveats but to
make explicit, for each threat, what was done to mitigate it and which
expansion item closes the residual gap.

\subsection{Construct Validity}
\label{sec:threats_construct}

Construct validity asks whether the measured quantity --- the marginal
accuracy range of a pipeline stage --- actually captures ``stage importance''
as claimed.

\begin{enumerate}
  \item \textbf{Range is confounded by the number of stage levels.} The
    headline comparison (7 vectorizers vs.\ 3--4 filtrations) stacks the deck:
    a stage with more levels can span a larger range by chance. \emph{Severity:
    moderate.} \emph{Mitigation (done):} the $\omega^2$ population effect size
    corrects for level counts (ECG200 vectorizer $\omega^2{=}0.165$, CI
    [0.063, 0.423], vs.\ filtration $-0.017$; MNIST vectorizer 0.214 vs.\
    filtration 0.143 --- the vectorizer CI excludes zero on both real datasets,
    the ECG200 filtration CI straddles it); the equal-footing analysis matches
    best-3 vectorizers against 3 filtrations; the two-way interaction ANOVA
    (\S4.1) shows the stage decomposition is not an artefact of
    main-effects-only modelling. \emph{Residual:} the levels-matched analysis
    still compares the \emph{best} three vectorizers, which is mildly
    favourable to the vectorizer stage.
  \item \textbf{Scalar vectorizers create a floor effect.} Persistence Entropy
    and Amplitude are single scalars; their poor accuracy inflates the
    vectorizer range from the bottom. \emph{Severity: moderate.}
    \emph{Mitigation (done):} ranges are also reported excluding the
    degenerate scalar vectorizers (ECG200 vectorizer range contracts from
    6.39pp to 3.10pp but remains the largest stage; \S4.1). \emph{Residual:}
    the exclusion rule was applied post hoc, so the headline 6.39pp includes
    the floor effect by design and the paper says so.
  \item \textbf{Single-split point estimates are noisy.} Per-configuration
    accuracy varies by 1.09pp on average across the five splits (max 3.11pp;
    50 of 112 configurations exceed 1pp); the headline 83.0\% configuration
    drops to 79.6\% at $r{=}5$ (rank 4). \emph{Severity: high for any
    single-split number.} \emph{Mitigation (done):} all headline claims rest
    on the r$=$25 repeated-CV protocol with Nadeau--Bengio corrected CIs
    (vectorizer [5.69, 7.10], classifier [2.92, 4.07], filtration [0.37,
    1.01], all excluding zero) and Friedman/sign-test inference on the
    repetition-level ordering; single-split numbers are labelled as such
    throughout. \emph{Residual:} MNIST binary is r$=$5, ECG5000 and the panel
    were single-split at the time of writing (expansion \#5 closes this).
  \item \textbf{Accuracy is the only outcome for most of the sweep.} F1,
    precision, and recall are stored per fold but the analysis is
    accuracy-centric. \emph{Severity: low-to-moderate.} \emph{Mitigation
    (done):} beyond-accuracy analysis on ECG5000 (balanced accuracy, per-class
    breakdowns; expansion \#14). \emph{Residual:} calibration and AUROC are
    not uniform across datasets.
\end{enumerate}

\subsection{Internal Validity}
\label{sec:threats_internal}

Internal validity concerns confounds and measurement artefacts within the
study.

\begin{enumerate}
  \item \textbf{Single implementation library.} Every filtration and
    vectorizer runs through giotto-tda. \emph{Severity: moderate} --- the
    weak-Alpha fragility (IndexError on quantized UCR series) proves
    implementation dependence is real. \emph{Mitigation (planned):}
    cross-library replication in GUDHI-native and Ripser-native paths
    (expansion \#11), with the fragility reported as a first-class result.
  \item \textbf{Homology capped at $H_1$.} $H_2$ requires $O(n^4)$ simplices
    for Vietoris--Rips; the torus's $\beta_2$ is never measured.
    \emph{Severity: moderate.} \emph{Mitigation (planned):} Alpha complex in
    3D for $H_2$ (expansion \#9).
  \item \textbf{Binary-only classification.} All datasets except ECG5000
    (5 classes) are binary. \emph{Severity: moderate} --- the 10-class MNIST
    probe shows the stage ordering flips at multi-class scale (filtration
    4.53pp vs.\ vectorizer 3.44pp), so the binary result is a boundary
    condition, not a law (\S4.2). \emph{Mitigation (done):} the flip is
    reported; \emph{planned:} multi-class breadth (expansion \#6 datasets).
  \item \textbf{Fixed hyperparameters.} No grid search on any stage, by
    design. \emph{Severity: low} --- the hyperparameter arm (B3) shows
    vectorizer dominance contracts but persists under one-parameter-at-a-time
    tuning (ECG200 5.75$\rightarrow$4.75pp; MNIST 1.75$\rightarrow$1.62pp), so
    dominance is not a default-settings artefact. \emph{Residual:} best-tuned
    values carry mild selection optimism (selected on the same folds that
    score the range), which would \emph{inflate} the tuned range --- the
    observed contraction is therefore conservative.
  \item \textbf{Subsampling choices.} Uniform random subsampling caps point
    clouds; farthest-point sampling was deferred. \emph{Severity: low.}
    \emph{Mitigation (done):} the FPS ablation (B4) finds no benefit
    ($-0.25$pp overall; uniform wins at $k{=}15$, $\sigma{=}0.30$).
  \item \textbf{The sphere/torus norm confound.} The classes differ in scale
    as well as topology. \emph{Severity: moderate.} \emph{Mitigation (done):}
    the matched genus-1/genus-2 control (identical norm distributions) shows
    TDA retains 95.83\% where norm features collapse to 48--58\% at
    $\sigma{=}0.30$. \emph{Residual caveat:} the matched pair is still
    linearly separable in raw coordinates (logistic 100\% at both noise
    levels, \texttt{baseline\_experiments.db}) --- the control isolates the
    norm/scale confound, not linear separability in general.
  \item \textbf{Determinism and exclusion discipline.} Seeds are fixed (base
    42; per-repetition 43--67; per-dataset CRC32 subsampling), and the
    56/672 excluded configurations (point-cloud filtrations on image data)
    are counted and disclosed rather than silently dropped. ECG200 is
    energy-normalised (sum of squares constant at 95.00 across all 200
    signals), so the trivial-separator concern does not apply to the executed
    data (\S4.1).
\end{enumerate}

\subsection{External Validity}
\label{sec:threats_external}

External validity asks how far the findings generalise beyond the studied
configuration space.

\begin{enumerate}
  \item \textbf{Dataset breadth.} Two time series (ECG200, ECG5000), one
    image pair (binary MNIST), and synthetic point clouds. \emph{Severity:
    high.} \emph{Mitigation (done):} the 9-dataset panel (B2) replicates the
    vectorizer $>$ filtration ordering on 7/9 datasets (median 3.32 vs.\
    1.37pp; filtration wins on HandOutlines and 10-class MNIST), and the
    within-filtration vectorizer span remains large even where filtration
    leads (10-class MNIST: 9--10pp). \emph{Planned:} topology-wins datasets
    (\#6), further UCR coverage, multi-patient ECG (MIT-BIH).
  \item \textbf{Single-patient ECG5000.} ECG5000 is BIDMC chf07 --- one
    patient. The time-series result could be patient-specific.
    \emph{Severity: moderate.} \emph{Mitigation (planned):} multi-patient
    MIT-BIH arrhythmia sweep; the panel adds ECG200/ECG5000 diversity but not
    patient diversity.
  \item \textbf{Hardware dependence of runtime results.} Wall-clock numbers
    come from one workstation (8-core, 16GB). \emph{Severity: low for
    ordering claims} (the 3--27\% filtration speed gap is structural: simplex
    counts differ by construction, \S4.3), \emph{moderate for absolute
    numbers}. \emph{Mitigation:} per-configuration wall times and peak memory
    are stored in the results DBs; peak memory was not measured per stage.
  \item \textbf{Missing method families.} No learned vectorizers (PersLay,
    Hofer; \#8), no $H_2$ features (\#9), no cross-library replication
    (\#11), no protein/graph modalities. \emph{Severity: moderate} --- each
    is a named boundary of the current claim, and each maps to a registered
    expansion. The paper's contribution is scoped accordingly: the internal
    stage-importance decomposition and the framework, not an
    absolute-accuracy win over classical methods (raw baselines beat TDA on
    the studied sets; \S4.3, \S5.2).
\end{enumerate}

%% END OF DRAFT — reorganises \S5.2's Limitations paragraph into the three
%% standard categories; safe to insert as-is or to replace the Limitations
%% paragraph once the author has merged the two.
